\section{Stretching rounds and the width bound}\label{sec:assembly}
\label{assembly:section}

We now construct the metric and log warps required by
Theorem~\ref{interface:width}.  The local stretching lemma increases
one separation at a time.  Its loss is measured by a decrease in the
trace of a fixed background metric, so a whole family of separations
can be treated without a bound on the overlap of their neighborhoods.
Changing the working scale after each round restores the quadratic
reserve.

\subsection{Summing the local losses}

The following form of the round argument accommodates a drift of
either sign.  It extends the iteration in
\cite[Section~4]{OpenAIWidth} using the positive-part term in
Lemma~\ref{stretch:local}.

\begin{lemma}[One round of separations]\label{assembly:round}
Let $n\ge4$, let
$0<\delta<\min\{1/10,1/(n+1)\}$, and let $W,T$ be supplied by
Lemma~\ref{stretch:local}.  Put $V(t)=t^2+T^2$ and $L_*=1.004$.
Let $h_0$ be a complete smooth metric on a connected smooth
$n$-manifold $M$ without boundary, and let $t_0:M\to\R$ be smooth.
Suppose that, for a real constant $c$,
\[
 \Dcal(h_0,t_0)\ge c+V(t_0).
\]
Let $(A_i,B_i)$ be a finite or countable family of nonempty closed
pairs, with each $A_i$ compact and
$\dist_{h_0}(A_i,B_i)\ge W$.  Suppose there is a locally finite family
of compact sets $K_i$ such that
\[
 \{x:\dist_{h_0}(A_i,x)<W\}\subset K_i\setminus B_i.
\]
Then there are smooth $h_1\ge h_0$ and $t_1\ge t_0$ such that
\begin{align}
 \dist_{h_1}(A_i,B_i)&\ge L_*W\quad\hbox{for every }i,
       \label{assembly:round-distance}\\
 \Dcal(h_1,t_1)&\ge c+\bigl(1-(n+1)\delta\bigr)V(t_1).
       \label{assembly:round-reserve}
\end{align}
The changes are a locally finite sequence of compactly supported
changes inside the sets $K_i$, and $h_1$ is complete.
\end{lemma}

\begin{proof}
Enumerate the pairs and choose $d_i>0$ with $\sum_i d_i<\delta$.
Apply Lemma~\ref{stretch:local} to them successively, keeping $h_0$
as the background metric throughout the round.  The separation
hypothesis is therefore fixed, and increasing the current metric
preserves its domination of $h_0$.  For the current metric and drift
$(H,z)$, define the pointwise functions
\[
 Q=\tr_H h_0,\qquad
 \Lambda=1-\delta(n-Q)-\Sigma,
\]
where $\Sigma$ is the sum of the losses $d_i$ already used.  We claim
that the inequality maintained after every finite number of steps is
\begin{equation}\label{assembly:trace-invariant}
 \Dcal(H,z)\ge c+\Lambda V(z).
\end{equation}
Initially $Q=n$, $\Sigma=0$, and $\Lambda=1$.  At every subsequent
step $0<Q\le n$ and $\Sigma<\delta$, hence
$0<\Lambda\le1$.  Denote the next metric, drift, trace, and coefficient
by $H',z',Q',\Lambda'$, and its small loss by $d$.  Then
$\Lambda'=\Lambda-\delta(Q-Q')-d$.  The local lemma and
\eqref{assembly:trace-invariant} give
\begin{align*}
 \Dcal(H',z')
 &\ge c+\Lambda V(z)+(V(z')-V(z))_+
                 -\bigl(\delta(Q-Q')+d\bigr)V(z')\\
 &\ge c+\Lambda'V(z').
\end{align*}
Here $(a)_+=\max\{a,0\}$, and the second inequality follows from
$a_+\ge\Lambda a$ for $0\le\Lambda\le1$.  This comparison is valid
even when $V(z')<V(z)$.  Since
$\Lambda\ge1-(n+1)\delta$, it proves the asserted reserve at every
finite stage.

Each processed pair attains separation at least $L_*W$, and later
metric enlargements preserve that bound.  The changes are supported
in the locally finite compact family $(K_i)$.  Thus near every point
the metric and drift are unchanged after finitely many steps.  Their
limits $h_1,t_1$ are smooth and satisfy the same local scalar inequality;
metric monotonicity preserves every separation in the limit.
Finally, a smooth metric dominating a complete metric is complete:
an $h_1$-Cauchy sequence converges in $h_0$, and local equivalence of
the two smooth metrics gives convergence in $h_1$.
\end{proof}

\subsection{Proof of Theorem~\ref{thm:main}}

\begin{proof}
\textbf{Initialization.}
Fix $n\ge4$.  Choose $\delta>0$, depending only on $n$, so that
\begin{equation}\label{assembly:delta}
 \delta<\frac1{10},\qquad
 2(n+1)\delta<1-L_*^{-2},\qquad L_*=1.004.
\end{equation}
Take $W,T$ from Lemma~\ref{stretch:local}, and write
$V(t)=t^2+T^2$.  Let $C(n)$ be the constant in
Proposition~\ref{pre:reserve}.  Choose $R_0\ge1$, depending only on $n$,
such that
\begin{equation}\label{assembly:scale}
 \frac{R_0^2}{2}\ge C(n)+T^2,
 \qquad WR_0\ge D_{\mathrm{nat}}(n),
 \qquad D=WR_0.
\end{equation}
These choices precede all metric-dependent constructions.

Let $(M,g)$ satisfy the spectral hypothesis; we may assume $M$ is
nonempty.  Choose a maximal
$D$-separated set of centers $(x_i)_{i\in I}$, and set
\[
 A_i=\overline B_g(x_i,D),\qquad
 B_i=M\setminus B_g(x_i,2D),\qquad
 K_i=\overline B_g(x_i,2D).
\]
The $A_i$ cover $M$, and properness of $g$ makes $(K_i)$ locally finite:
centers whose outer balls meet a compact set lie in a compact
$2D$-neighborhood of that set, which contains only finitely many
$D$-separated points.  Second countability makes $I$ finite or
countable.  If some $B_i$ is empty, then $\diam_gM\le4D$ and a
constant map proves the theorem.  Henceforth every $B_i$ is nonempty.
Choose positive numbers $p_i^*$, including $p_0^*$, with total $1/20$,
and prescribe
\begin{equation}\label{assembly:targets}
 L_i=10+2\pi\sqrt{2/p_i^*}\quad(i\in I).
\end{equation}

By Lemma~\ref{pre:supersolution}, there is a smooth real function $w$
with $\Dcal(g,w)\ge1$.  Dividing a metric by $R^2$ multiplies scalar
curvature, the Laplacian, and squared gradient norms by $R^2$.
Consequently
$\Dcal(g/R_0^2,w)\ge R_0^2$.  Apply Proposition~\ref{pre:reserve} to
$(g/R_0^2,w)$, and multiply its output metric by $R_0^2$.
Call the resulting metric $h^{[0]}$ and its drift $t^{[0]}$.
Then $h^{[0]}\ge g$ is complete and
\begin{align}
 \Dcal(h^{[0]}/R_0^2,t^{[0]})
 &\ge R_0^2+(t^{[0]})^2-C(n)\notag\\
 &\ge \frac{R_0^2}{2}+V(t^{[0]}).
       \label{assembly:initial-reserve}
\end{align}
The original drift $w$ may be unbounded.  The uniform scale in
\eqref{assembly:scale} is available because preprocessing supplies
the quadratic term in the new drift.

\medskip\noindent\textbf{Successive rounds.}

Put $R_j=L_*^jR_0$.  At the start of round $j$ we maintain complete
smooth data $(h^{[j]},t^{[j]})$ satisfying
\begin{equation}\label{assembly:start}
 \Dcal(h^{[j]}/R_j^2,t^{[j]})
       \ge\frac{R_j^2}{2}+V(t^{[j]}).
\end{equation}
A label is called active when
$\dist_{h^{[j]}}(A_i,B_i)\le L_i$; every active label will also satisfy
\begin{equation}\label{assembly:active-distance}
 \dist_{h^{[j]}}(A_i,B_i)\ge WR_j.
\end{equation}
Both assertions hold at $j=0$: the reserve is
\eqref{assembly:initial-reserve}, and $h^{[0]}\ge g$ gives separation
at least $D=WR_0$.

Use the working metric $h_0=h^{[j]}/R_j^2$ for round $j$.
For each active label, its open $h_0$-distance-$W$ neighborhood of
$A_i$ misses $B_i$, by \eqref{assembly:active-distance}.  Therefore
it lies in the fixed original ball $B_g(x_i,2D)\subset K_i$.
Lemma~\ref{assembly:round}, with $c=R_j^2/2$, applies to the active
pairs.  Multiply its output metric by $R_j^2$ to obtain $h^{[j+1]}$,
and denote its output drift by $t^{[j+1]}$.  If there are no active
pairs, keep the metric and drift unchanged.  In either case,
$h^{[j+1]}\ge h^{[j]}$.  Every label still active at the next round
now has separation at least $L_*WR_j=WR_{j+1}$.

The reserve is restored by the new working scale:
\begin{align*}
 \Dcal(h^{[j+1]}/R_{j+1}^2,t^{[j+1]})
 &\ge\frac{R_{j+1}^2}{2}
       +L_*^2\bigl(1-(n+1)\delta\bigr)V(t^{[j+1]})\\
 &\ge\frac{R_{j+1}^2}{2}+V(t^{[j+1]}),
\end{align*}
where \eqref{assembly:delta} makes the coefficient of $V$ greater
than one.  Thus \eqref{assembly:start} holds at every round.  In
original units its constant scalar reserve remains $1/2$.

\medskip\noindent\textbf{The limit and its log warps.}

Each label is active for only finitely many rounds, because its
separation is at least $WR_j$ while active and $R_j\to\infty$.
It remains inactive once its target is exceeded, since all metrics
increase.  Every change associated with label $i$, in every round,
is supported in the same original compact ball $K_i$.  A compact
set meets only finitely many such balls, and each of their labels
is used only finitely many times.  Hence on a neighborhood of any
compact set the metric and drift are unchanged after finitely many
changes in the entire construction.  We obtain smooth data $(h,t)$
with
\begin{equation}\label{assembly:final-drift}
 h\ge g,\qquad \Dcal(h,t)\ge\frac12,
 \qquad \dist_h(A_i,B_i)>L_i\quad(i\in I).
\end{equation}
The metric $h$ is complete by domination of $g$.

It remains to realize the drift by log warps while spending at most
$1/10$ of the scalar reserve.  Dividing a log warp into many equal
parts preserves its sum and reduces the sum of its squared gradients;
see \cite[Section~1.A]{Gromov2024}.  We use the localized construction
of \cite[Theorem~4.3]{OpenAIWidth}, applied to a compactly supported
decomposition of our signed drift.  Take a finite or countable smooth
partition of unity $(\varphi_\nu)_{\nu\ge1}$ with locally finite
compact supports, and put $\tau_\nu=-\varphi_\nu t$.  Then
$\sum_\nu\tau_\nu=-t$.  For each $\nu$, choose an integer $m_\nu\ge1$
so large that
\[
 \frac{1}{m_\nu}\sup_M|\dd\tau_\nu|_h^2
          \le\frac{2^{-\nu}}{10}.
\]
The supremum is finite because $\tau_\nu$ has compact support.
Introduce $m_\nu$ log warps all equal to $\tau_\nu/m_\nu$, and list
all these functions as $(p_\alpha)$.  Repeating each support only
finitely many times preserves local finiteness.  Therefore
\[
 P=\sum_\alpha p_\alpha=-t,
 \qquad
 \sum_\alpha|\dd p_\alpha|_h^2
   =\sum_\nu\frac{|\dd\tau_\nu|_h^2}{m_\nu}\le\frac1{10}.
\]
Substitution into \eqref{interface:scalar} gives
\[
 R(h,(p_\alpha))
 =\Dcal(h,t)-\sum_\alpha|\dd p_\alpha|_h^2
 \ge\frac25.
\]

Equations \eqref{assembly:targets} and \eqref{assembly:final-drift}
now supply all hypotheses of Theorem~\ref{interface:width}.  Its map
has whole fibers of diameter at most $4D+2$ in the original metric
$g$.  Taking
\[
 C_n=4WR_0+2
\]
proves the theorem, since $W$ and $R_0$ depend only on $n$.
\end{proof}

For a quadratic-form lower bound $\lambda>0$, the metric $\lambda g$
has normalized lower bound $1$, since its scalar curvature and
Laplacian are $\lambda^{-1}$ times those of $g$.  Its distances are
$\sqrt\lambda$ times those of $g$, giving the bound
$C_n/\sqrt\lambda$ stated in the introduction.
